% ============================================================
% บรรณานุกรม (References - APA 7th Edition)
% ============================================================
\chapter*{บรรณานุกรม}
\addcontentsline{toc}{chapter}{บรรณานุกรม}

\begin{hangparas}{1.25cm}{1}

Connectivity Standards Alliance. (2023). \textit{Zigbee specification: Document 05-3474-22}. CSA-IOT. \url{https://csa-iot.org/}

Espressif Systems. (2024). \textit{ESP32-C6 technical reference manual: 2.4 GHz Wi-Fi 6, Bluetooth 5, and IEEE 802.15.4 SoC}. Espressif Systems Co., Ltd. \url{https://www.espressif.com/}

IEEE Standards Association. (2020). \textit{IEEE standard for low-rate wireless networks: IEEE Std 802.15.4-2020}. IEEE. \url{https://doi.org/10.1109/IEEESTD.2020.9144691}

Koen, K. (2024). \textit{Zigbee2MQTT documentation and device ecosystem}. Zigbee2MQTT Open Source Project. \url{https://www.zigbee2mqtt.io/}

Modbus Organization. (2022). \textit{Modbus application protocol specification V1.1b3}. Modbus-IDA. \url{https://modbus.org/}

Texas Instruments. (2023). \textit{CC2652P SimpleLink multi-protocol 2.4 GHz wireless MCU with integrated power amplifier}. Texas Instruments Incorporated. \url{https://www.ti.com/}

ชีวะ ทัศนา, จิรภัทร จันทมาลี, อรรถกร คำฉัตร, และ ณมนรัก คำฉัตร. (2568). สถาปัตยกรรมโครงข่ายเซนเซอร์ไร้สายและปัญญาประดิษฐ์ฝังตัวสำหรับการตรวจวัดความชื้นดินและบรรยากาศในแปลงผลไม้มูลค่าสูงภาคตะวันออก. \textit{วารสารวิทยาศาสตร์และเทคโนโลยี มหาวิทยาลัยราชภัฏรำไพพรรณี}, \textit{9}(2), 45--58.

ชีวะ ทัศนา. (2567). การวิเคราะห์และออกแบบระบบไอโอทีสำหรับการเกษตรแม่นยำสูงในสวนทุเรียน. ใน \textit{การประชุมวิชาการระดับชาติด้านวิทยาศาสตร์ประยุกต์และวิศวกรรมศาสตร์} (หน้า 112--125). มหาวิทยาลัยเชียงใหม่.

นิภัทร เปี่ยมอรุณ, นันทพร มูลรังษี, และ ชีวะ ทัศนา. (2568). แบบจำลองการแพร่กระจายคลื่นวิทยุความถี่ 2.4 GHz ในทรงพุ่มสวนผลไม้เขตร้อนชื้น. \textit{วารสารวิจัยและพัฒนานวัตกรรมดิจิทัล}, \textit{5}(1), 22--35.

\end{hangparas}

\clearpage
